\chapter{Ferromagnetic random variables}\label{chap:ferro}

A symmetric random variable is ferromagnetic if its moment generating function vanishes only on the imaginary axis. By \eqref{eq:ch2:mgf} the moment generating function of the P\'olya variable $X$ is $\xi(\frac12+s)/\xi(\frac12)$, so RH is the statement that one explicit law, the P\'olya law with density $p=\Phi/\Xi(0)$, is ferromagnetic.

With no hypothesis, the P\'olya law lies within total variation distance $1.3\times10^{-4}$ of a law that is ferromagnetic by an explicit construction, and within $2\times10^{-6}$ of one whose positivity, and hence ferromagnetism, rests on a computer assisted proof in ball arithmetic (Section~\ref{ch:fmapprox}).

The scripts are in the folder \texttt{ferro} (\texttt{code/\allowbreak README.md}). We use the units of \eqref{eq:ch2:Phi} throughout, in which the zeros of the characteristic function of $X$ are the ordinates $\gamma$ and $\operatorname{Var}X=2\E T=0.0462099862$. The normalisation of Rodgers and Tao in Section~\ref{ch:dbn} doubles the variable of the kernel, and statements quoted from the literature have been converted. As in Section~\ref{ch:clock}, $Z$ is a standard normal variable, the $e_k$ are standard exponential variables, the $L_k$ are standard Laplace variables, $\varepsilon$ and the $\varepsilon_k$ are fair random signs, $U$ is uniform on $[-1,1]$, and all of them are independent of each other and of everything else.
